\section{The \texttt{multi\_pipe} McStas Component}
multi-pipe circular slit.

\subsection*{Identification}
\begin{itemize}
  \item \textbf{Site:} 
  \item \textbf{Author:} Uwe Filges
  \item \textbf{Origin:} PSI
  \item \textbf{Date:} March, 2005
\end{itemize}

\subsection*{Description}
\begin{lstlisting}
No transmission around the slit is allowed.

example for an input file
# user defined geometry
# x(i) y(i) r(i) [m]
0.02 0.03 0.01
-0.04 0.015 0.005

warning: at least two values must be in the file

Example: multi_pipe(xmin=-0.01, xmax=0.01, ymin=-0.01, ymax=0.01)
\end{lstlisting}

\subsection*{Input parameters}
Parameters in \textbf{boldface} are required; the others are optional.

\begin{longtable}{p{0.22\textwidth}p{0.12\textwidth}p{0.46\textwidth}p{0.14\textwidth}}
\toprule
\textbf{Name} & \textbf{Unit} & \textbf{Description} & \textbf{Default} \\
\midrule
\endhead
filename & str & define user table of holes & 0 \\
\textbf{xmin} & m & Lower x bound &  \\
\textbf{xmax} & m & Upper x bound &  \\
\textbf{ymin} & m & Lower y bound &  \\
\textbf{ymax} & m & Upper y bound &  \\
radius &  & radius of a single hole & 0.0 \\
gap & m & distance between holes & 0.0 \\
thickness & m & thickness of the pipe & 0.0 \\
xwidth & m & Width of slit plate. Overrides xmin,xmax. & 0 \\
yheight & m & Height of slit plate. Overrides ymin,ymax. & 0 \\
\bottomrule
\end{longtable}

\subsection*{Links}
\begin{itemize}
  \item \href{run:/home/willend/willend-McCode/mcstas-comps/contrib/multi_pipe.comp}{Source code} for \texttt{multi\_pipe.comp}.
\end{itemize}
\IfFileExists{multi_pipe_static.tex}{\input{multi_pipe_static.tex}}{}